\section*{Övningar}
\addcontentsline{toc}{section}{Övningar}

Lösningarna till uppgifterna som inte ber om kod står i \hyperref[ch:losningar]{appendix~A}.
Koden publiceras i kursrepot efter föreläsningen.

\exercise{Bygg en frame}{Räkning}
\label{ex:1:1}
En nod med adress \code{0x17} skickar en \code{StatusRequest} till en temperatursensor med adress
\code{0x25}. Sekvensnumret är \code{0x7F05}.
\begin{enumerate}[a)]
  \item Skriv upp hela framen byte för byte, på hexadecimal form.
  \item Sensorn svarar med ett \code{StatusResponse} som innehåller temperaturen \code{0x16}
    (en byte). Skriv upp den framen också.
  \item Hur många bytes är den kortast möjliga framen i det här protokollet?
\end{enumerate}

\exercise{Fälten och deras syfte}{Resonemang}
\label{ex:1:2}
För vart och ett av fälten SOF, LEN, TYPE, DST, SRC, SEQ och CHK: beskriv i en mening vad som
går sönder om fältet tas bort.

Två av dem går att ta bort utan att protokollet slutar fungera i ett system med exakt två noder.
Vilka, och varför?

\exercise{Endianness}{Resonemang}
\label{ex:1:3}
En grupp serialiserar sekvensnumret genom att kopiera minnet direkt:

\begin{cppcode}
std::memcpy(&buf[Offset::Seq], &frame.seqNr, sizeof(frame.seqNr));
\end{cppcode}

\begin{enumerate}[a)]
  \item Vad hamnar i bufferten på en x86-maskin om \code{seqNr} är \code{0x7F05}?
  \item Koden fungerar ändå felfritt i deras egna tester. Varför?
  \item När går den sönder?
\end{enumerate}

\exercise{Checksummans gränser}{Resonemang}
\label{ex:1:4}
Kursens checksumma är en summa av bytes.
\begin{enumerate}[a)]
  \item En störning byter plats på två databytes. Upptäcks det?
  \item En störning ökar en byte med 1 och minskar en annan med 1. Upptäcks det?
  \item Ge ett exempel på en ändring i framen ovan som checksumman \emph{inte} upptäcker.
\end{enumerate}
